\section{Related Work}
\label{sec:related-work}

\draftnote{DRAFT. Carried over from an earlier draft on this topic and
rewritten in neutral third person; citations verified or sourced per
main.bib's own annotations (VERIFIED vs.\ SOURCED). Needs integration
with the overlearning/adversarial-leakage citations already used in
the adaptive-attacker appendix (Elazar \& Goldberg, Song \& Shmatikov)
-- those belong here too, not only in the appendix.}

\textbf{Fairness in machine learning.} Fairness-aware learning spans
pre-processing, in-processing, and post-processing
approaches~\cite{hardt2016equality}, and is commonly split into
\emph{group} fairness, which targets parity across demographic
groups, and \emph{individual} fairness, which requires similar
individuals to receive similar
treatment~\cite{dwork2012fairness,verma2018fairness}. Demographic
attributes are typically absent from anonymized trajectory datasets,
making individual fairness the natural fit for this setting: here, it
requires that users with similar mobility behavior experience
comparable privacy protection after anonymization, not that any
demographic group be treated comparably to another. Certifying and
correcting disparate impact once introduced is itself an active area
of study~\cite{feldman2015certifying,barocas2016big,mehrabi2021survey}.

\textbf{Privacy-preserving spatial-temporal models.} Classical
location-privacy mechanisms such as $k$-anonymity~\cite{sweeney2002k}
and geo-indistinguishability~\cite{andres2013geo} reduce identifiability
by perturbing or generalizing released locations directly, independent
of any learned model. More recent privacy-utility trade-off (PUT)
models instead learn a representation of mobility data that preserves
downstream utility while suppressing identifying information.
TrajGAN~\cite{rao2020lstm} trains a generative model to produce
synthetic trajectories that retain the spatial, temporal, and semantic
structure of real mobility traces while protecting the identity of the
individuals they are drawn from; this paper's baseline comparison
(Appendix~\ref{sec:adaptive-appendix}'s companion positive-control
analysis) ports and validates this same architecture. Mo-PAE instead
learns an encoder via adversarial multi-task training, trained to
retain information relevant to next-location prediction while
suppressing information relevant to re-identifying the
individual~\cite{zhan2023mopae}. A direct empirical audit of models in
this family found that they can pass group-fairness checks while
violating individual fairness, producing disparate privacy and utility
outcomes for users with otherwise similar mobility behavior, when
evaluated with the in-model adversary's own naive, pairwise,
threshold-based metric~\cite{zhan2023analysing}. This paper
investigates that same question with a hardened, external evaluation
protocol (Section~\ref{sec:evaluation-protocol}) rather than the
in-model metric the original audit relied on, and finds the picture is
more specific, and in one respect different, than that audit's
own measurement could show (Section~\ref{sec:methods-changed}).

\textbf{Fairness in mobility and transportation.} Equity analysis in
transportation research has traditionally relied on coarser,
region- or group-level disparity measures, such as a region-based
fairness gap in ride-sharing
demand~\cite{yan2020fairness,yan2019fairst} or a systematic framework
for benchmarking fairness in generative crowd-flow
models~\cite{he2025framework}. These measures are well suited to
aggregate planning questions but not to an individual-level claim
about who, specifically, is protected by a given privacy mechanism --
the question this paper asks directly.

\textbf{Representation leakage beyond the intended adversary.} Separately
from the mobility-privacy literature, work studying adversarially
trained representations in other domains has repeatedly found that an
in-training adversary's apparent success is not evidence that the
information it was trained against is actually absent from the
representation: a differently targeted, post-hoc classifier can
recover demographic attributes an adversarial removal method was
designed to suppress~\cite{elazar2018adversarial}, and models trained
on one narrow objective can be shown to ``overlearn'' and encode
substantially more than that objective required, including attributes
never labeled during
training~\cite{song2020overlearning}. This paper's own in-model-vs.-external
attacker comparison (Section~\ref{sec:residual-risk}) and gap
decomposition (Appendix~\ref{sec:adaptive-appendix}) are a mobility-privacy
instance of the same underlying phenomenon.
